\section*{Problem 3}



\partlabel{(a)}
% L13 notes PDF pp.1--2, Eqs. (267)--(281): Kantrowitz criterion.
% L9 notes PDF p.4, Eq. (27): normal-shock stagnation-pressure ratio.
Shock at capture plane: $M_{1s}=M_\infty$, $\dot m_1=\dot m_2$, $T_{0,1}=T_{0,2}$.
\begin{gather*}
\frac{A_c}{A_1^*}=F(M_\infty),\qquad
\frac{A_2^*}{A_1^*}=\frac{p_{0,1}}{p_{0,2}}=\frac{1}{\Pi(M_\infty)},\\
A_t\ge A_2^*
\quad\Longrightarrow\quad
\frac{A_t}{A_c}\ge\frac{1}{F(M_\infty)\Pi(M_\infty)}.
\end{gather*}
\begin{align*}
F(3)&=4.23456790,\qquad \Pi(3)=0.328343888,\\
\left(\frac{A_t}{A_c}\right)_{\min}
&=\frac{1}{(4.23456790)(0.328343888)}
=\boxed{0.719220342}<0.720000.
\end{align*}
$\boxed{\text{Self-starting at }M_\infty=3.}$

\partlabel{(b)}
% L12 notes PDF p.1, Eq. (228): area--Mach relation, supersonic branch.
% L13 notes PDF pp.3--4, Eq. (292) and following procedure: subsonic exit matching.
Shock at throat: $M_{1s}=M_t>1$.
% Bisection: throat [1,3]; subsonic exit [1e-7,1].
\begin{align*}
F(M_t)&=\frac{A_t}{A_1^*}
=\frac{A_t}{A_c}F(3)=0.72(4.23456790)=3.04888889,\\
M_t&=2.65452303,\qquad \Pi(M_t)=0.439907248,\\
p_{0,2}&=p_{0,\infty}\Pi(M_t)
=500(0.439907248)=\boxed{219.954~\mathrm{kPa}},\\
\Delta p_{0,\mathrm{loss}}&=p_{0,\infty}-p_{0,2}
=500-219.953624=\boxed{280.046~\mathrm{kPa}},\\
\frac{\Delta p_{0,\mathrm{loss}}}{p_{0,\infty}}
&=1-0.439907248=\boxed{56.0093\%}.
\end{align*}
% L12 notes PDF p.4, Eqs. (244)--(246); PDF p.2, Eq. (237).
$0<M_e<1$:
\begin{align*}
F(M_e)&=\frac{A_e}{A_2^*}
=\frac{A_e}{A_t}\frac{A_t}{A_1^*}\frac{A_1^*}{A_2^*}
=1.50(3.04888889)(0.439907248)\\
&=2.01184248
\quad\Longrightarrow\quad M_e=0.303882446,\\
p_{b,\max}&\simeq p_e
=\frac{p_{0,2}}{(1+0.2M_e^2)^{3.5}}
=\frac{219.953624}{[1+0.2(0.303882446)^2]^{3.5}}
=\boxed{206.307~\mathrm{kPa}}.
\end{align*}

\partlabel{(c)}
% L13 notes PDF p.2, Eqs. (271), (282): limiting starting equation.
% Bisection for the limiting starting Mach number: [3,8].
% L12 notes PDF p.1, Eq. (228), applied on the started supersonic branch.
% L9 notes PDF p.4, Eq. (27).
\begin{align*}
\text{(i)}\quad&
\frac{1}{F(M_{\infty,\mathrm{start}})\Pi(M_{\infty,\mathrm{start}})}=0.68
\quad\Longrightarrow\quad
\boxed{M_{\infty,\mathrm{start}}=3.78120}.\\[5pt]
\text{(ii)}\quad&
M_\infty=3:\quad F(M_t)=0.68F(3)=2.87950617
\quad\xrightarrow{M_t>1}\quad\boxed{M_t=2.59394}.\\[5pt]
\text{(iii)}\quad&
\frac{p_{0,2}}{p_{0,\infty}}=\Pi(2.59394374)
=\boxed{0.462412=46.2412\%},\quad
p_{0,2}=231.206~\mathrm{kPa}.\\[5pt]
\text{(iv)}\quad&
\Delta\!\left(\frac{p_{0,2}}{p_{0,\infty}}\right)
=0.462412096-0.439907248
=\boxed{2.25048\text{ percentage points}},\\
&\text{Relative improvement}
=\left(\frac{0.462412096}{0.439907248}-1\right)100\%
=\boxed{5.11582\%}.
\end{align*}

% Numerical fzero brackets: F(Mt), [1,3]; F(Me), [1e-7,1];
% starting equation, [3,8]. Standalone hw3_problem3.m documents the solves.
